\documentclass[12pt]{article}
\usepackage[utf8]{inputenc}
\usepackage[T1]{fontenc}
\usepackage[margin=1in]{geometry}
\usepackage{setspace}
\usepackage{xcolor}
\usepackage[colorlinks=true, urlcolor=blue!60!black, linkcolor=blue!60!black]{hyperref}
\setlength{\parindent}{2em}
\setlength{\parskip}{1.25em}
\renewcommand{\baselinestretch}{1.0}
\begin{document}
\begin{center}
{\LARGE \textit{Hicklin v. Orbeck}, 437 U.S. 518 (1978)\par}
\vspace{0.8em}
{\large Case Brief\par}
\end{center}
\vspace{1em}

\section{Memory Jogger}

Under the Article IV Privileges and Immunities Clause, a State may not leverage its ownership of natural resources to force private employers throughout the economic ripple effect of resource development to give a flat hiring preference to all state residents over nonresidents when nonresidents are not a ``peculiar source of the evil'' of local unemployment and the preference is not closely tailored to aiding the unemployed.

\section{Detailed Case Facts}

In 1972, citing ``uniquely high unemployment'' in the State, the Alaska Legislature enacted a statute titled ``Local Hire Under State Leases,'' commonly known as ``Alaska Hire,'' codified at Alaska Stat.\ Ann.\ \S\S\ 38.40.010--38.40.090 (1977) (\href{https://supreme.justia.com/cases/federal/us/437/518/}{Justia Opinion}; \href{https://www.courtlistener.com/opinion/109911/hicklin-v-orbeck/}{CourtListener}). Under \S\ 38.40.030(a), the Alaska Commissioner of Natural Resources was required to incorporate into ``all oil and gas leases, easements or right-of-way permits for oil or gas pipeline purposes, unitization agreements, or any renegotiation of any of the preceding to which the state is a party'' a provision mandating the employment of ``qualified Alaska residents'' in preference to nonresidents (\href{https://supreme.justia.com/cases/federal/us/437/518/}{437 U.S.\ at 520 \& n.2}). Implementing regulations further mandated that all nonresidents be laid off before any resident working in the same trade or craft could be terminated (8 Alaska Admin.\ Code 35.011 (1977)) and administered the preference through state-issued ``resident cards'' (8 Alaska Admin.\ Code 35.015 (1977)) (\href{https://supreme.justia.com/cases/federal/us/437/518/}{437 U.S.\ at 520 \& n.1}).

The Act's statutory reach, defined in \S\ 38.40.050(a), was sweeping: it covered ``all employment which is a result of'' covered state oil and gas leases, pipeline permits, or unitization agreements taking place inside Alaska, extending not only to direct lessees on state land, but also to contractors, subcontractors, and suppliers of goods or services to contractors or subcontractors, as well as downstream refineries and distribution systems using state-leased oil and gas (\href{https://supreme.justia.com/cases/federal/us/437/518/}{437 U.S.\ at 529--31 \& nn.13--14}). Under \S\ 38.40.090(1), a ``resident'' had to be physically present in Alaska for one continuous year, maintain a place of residence in the State, establish voting residency, claim no residency in another State, and demonstrate intent to make Alaska a permanent home (\href{https://supreme.justia.com/cases/federal/us/437/518/}{437 U.S.\ at 521--22 n.5}).

Although enacted in 1972, Alaska Hire was not seriously enforced until 1975, when construction of the Trans-Alaska Pipeline reached its peak (\href{https://supreme.justia.com/cases/federal/us/437/518/}{437 U.S.\ at 521}). On March 1, 1976, after receiving complaints that nonresidents were being dispatched to pipeline jobs ahead of qualified Alaskans, Alaska Commissioner of Labor Edmund Orbeck issued Executive Order \#76-1, a cease-and-desist order directing all unions supplying pipeline workers ``to respond to all open job calls by dispatching all qualified Alaska residents before any nonresidents are dispatched'' (\href{https://supreme.justia.com/cases/federal/us/437/518/}{437 U.S.\ at 521}). Eight individuals---Sidney S.\ Hicklin, Ruby E.\ Dorman, Betty Cloud, Tommy Ray Woodruff, Frederick A.\ Mathers, Harry A.\ Browning, Emmett Ray, and Joseph G.\ O'Brien---all but one of whom had previously worked on the pipeline, were barred from obtaining pipeline-related jobs because they could not obtain resident cards (\href{https://supreme.justia.com/cases/federal/us/437/518/}{437 U.S.\ at 521, 523 \& n.7}; \href{https://law.justia.com/cases/alaska/supreme-court/1977/3025-1.html}{565 P.2d 159 (Alaska 1977)}).

\section{Procedural History}

On April 28, 1976, the eight plaintiffs filed suit in the Superior Court of Alaska in Anchorage against Labor Commissioner Edmund Orbeck and Natural Resources Commissioner Guy Richard Martin, seeking declaratory and injunctive relief on the grounds that Alaska Hire violated both the Privileges and Immunities Clause of Article IV, \S\ 2, and the Equal Protection Clause of the Fourteenth Amendment (\href{https://supreme.justia.com/cases/federal/us/437/518/}{437 U.S.\ at 520--22}; \href{https://law.justia.com/cases/alaska/supreme-court/1977/3025-1.html}{565 P.2d 159}). By agreement of the parties, the preliminary injunction motion was consolidated with a decision on the merits, submitted on affidavits, depositions, and legal memoranda without oral testimony. On July 21, 1976, the Superior Court upheld Alaska Hire in its entirety and denied all relief (\href{https://supreme.justia.com/cases/federal/us/437/518/}{437 U.S.\ at 522}).

On June 3, 1977, the Alaska Supreme Court unanimously invalidated Alaska Hire's one-year durational residency requirement under both the state and federal Equal Protection Clauses, holding further that any durational residency requirement exceeding 30 days was constitutionally infirm (\href{https://law.justia.com/cases/alaska/supreme-court/1977/3025-1.html}{\textit{Hicklin v. Orbeck}, 565 P.2d 159, 165, 171 (Alaska 1977)}). By a 3--2 vote, however, the Alaska Supreme Court upheld the Act's general preference for bona fide Alaska residents over nonresidents, relying on Alaska's ownership of its oil and gas resources under \textit{McCready v. Virginia}, 94 U.S.\ 391 (1877) (\href{https://supreme.justia.com/cases/federal/us/437/518/}{437 U.S.\ at 522, 528}). The plaintiffs appealed the 3--2 affirmance of the general resident hiring preference to the U.S.\ Supreme Court (the State did not cross-appeal the invalidation of the one-year requirement). The Supreme Court noted probable jurisdiction (434 U.S.\ 919 (1977)), heard oral argument on March 21, 1978, and issued its decision reversing the Alaska Supreme Court on June 22, 1978 (\href{https://supreme.justia.com/cases/federal/us/437/518/}{437 U.S.\ at 518, 522}; \href{https://www.oyez.org/cases/1977/77-324}{Oyez}).

\section{Judicial Votes}

The Supreme Court decided the case \textbf{9--0} (unanimous) (\href{https://supreme.justia.com/cases/federal/us/437/518/}{437 U.S.\ at 519}; \href{https://www.oyez.org/cases/1977/77-324}{Oyez}):
\begin{itemize}
    \item \textbf{Majority Opinion (9):} Justice William J.\ Brennan, Jr.\ wrote the opinion of the Court, joined in full by Chief Justice Warren E.\ Burger and Justices Potter Stewart, Byron R.\ White, Thurgood Marshall, Harry A.\ Blackmun, Lewis F.\ Powell, Jr., William H.\ Rehnquist, and John Paul Stevens.
    \item \textbf{Concurring Opinions (0):} None.
    \item \textbf{Dissenting Opinions (0):} None.
\end{itemize}

\section{Holding}

\textbf{Judgment:} Reversed (\href{https://supreme.justia.com/cases/federal/us/437/518/}{437 U.S.\ at 519, 534}). Alaska Hire violates the Privileges and Immunities Clause of Article IV, \S\ 2. (The Court also held as a threshold matter that the case was not moot: although the claims of the three short-term Alaska resident appellants---Hicklin, Dorman, and Browning---appeared moot once the one-year durational requirement fell, at least five nonresident appellants---Woodruff, Mathers, Ray, Cloud, and O'Brien---still could not qualify as residents under \S\ 38.40.090(1)(D) and retained a live controversy (\href{https://supreme.justia.com/cases/federal/us/437/518/}{437 U.S.\ at 523 \& n.7}).)

\textbf{Rule Announced:} Under the Article IV Privileges and Immunities Clause (where ``citizen'' and ``resident'' are essentially interchangeable for most analytical purposes), a State may not discriminate against out-of-state citizens in pursuing a common calling or private employment unless (1) there is a ``substantial reason'' for the discrimination---specifically, a showing that nonresidents constitute a ``peculiar source of the evil'' the statute seeks to remedy---and (2) the discrimination bears a ``substantial relationship'' (close tailoring) to that particular evil (\href{https://supreme.justia.com/cases/federal/us/437/518/}{437 U.S.\ at 524--28}). Furthermore, a State's ownership of a natural resource does not create a categorical exemption from the Clause; it is merely a factor---albeit often a crucial one---that cannot justify pervasive hiring discrimination forced upon private employers whose activities are attenuated from the State's proprietary interest, especially when the resource is bound for interstate commerce (\href{https://supreme.justia.com/cases/federal/us/437/518/}{437 U.S.\ at 528--34}).

\textbf{Holding vs.\ Dicta:}
\begin{itemize}
    \item \textbf{Binding Holding:} Alaska Hire's across-the-board resident hiring preference violated Article IV, \S\ 2 because (a) the record failed to show nonresidents were a peculiar source of Alaska's unemployment, (b) the flat preference for all Alaskans (employed or unemployed, skilled or unskilled) was not substantially related to alleviating unemployment, and (c) Alaska's ownership of oil and gas could not justify sweeping private-sector discrimination extending to contractors, subcontractors, suppliers, and downstream operations with little or no direct proprietary tie to the State.
    \item \textbf{Dicta / Reserved Questions:} The Court expressly declined to reach appellants' Fourteenth Amendment Equal Protection Clause challenge (\href{https://supreme.justia.com/cases/federal/us/437/518/}{437 U.S.\ at 534 n.19}). It also left open---while expressing doubt---whether a State may \textit{ever} require private employers to discriminate against nonresidents to combat local unemployment, and whether a more narrowly tailored preference limited solely to unemployed residents or residents enrolled in job-training programs might survive scrutiny (\href{https://supreme.justia.com/cases/federal/us/437/518/}{437 U.S.\ at 526, 528}). Finally, the Court's discussion of Commerce Clause natural-resource cases was explicitly offered as ``support'' for its Privileges and Immunities conclusion rather than as an independent Commerce Clause holding, because appellants raised no Commerce Clause claim (\href{https://supreme.justia.com/cases/federal/us/437/518/}{437 U.S.\ at 531--34}).
\end{itemize}

\section{Analysis of Opinions}

\subsection{Majority Opinion (Justice Brennan)}

Justice Brennan's unanimous opinion proceeded in four structured steps (\href{https://supreme.justia.com/cases/federal/us/437/518/}{437 U.S.\ at 523--34}; \href{https://www.courtlistener.com/opinion/109911/hicklin-v-orbeck/}{CourtListener}):

First, Justice Brennan framed the purpose of the Privileges and Immunities Clause of Article IV, \S\ 2, as establishing a ``norm of comity'' designed to ``constitute the citizens of the United States one people'' (quoting \textit{Paul v. Virginia}, 75 U.S.\ (8 Wall.) 168, 180 (1869)) and protecting the right of a citizen of one State to travel to another to ply a trade, practice an occupation, or pursue a common calling (\textit{Ward v. Maryland}, 79 U.S.\ (12 Wall.) 418 (1871)) (\href{https://supreme.justia.com/cases/federal/us/437/518/}{437 U.S.\ at 523--25}). Adopting the two-part analytical framework of \textit{Toomer v. Witsell}, 334 U.S.\ 385 (1948), and \textit{Mullaney v. Anderson}, 342 U.S.\ 415 (1952), the Court required both a ``substantial reason'' showing that nonresidents are a ``peculiar source of the evil'' targeted by the law and a ``reasonable'' or ``substantial'' relationship between the danger posed by nonresidents and the discrimination practiced upon them (\href{https://supreme.justia.com/cases/federal/us/437/518/}{437 U.S.\ at 525--27}).

Second, applying \textit{Toomer}, the Court found Alaska Hire failed both prongs. On the record---including a 1971 Federal Field Committee report cited by the State itself---Alaska's high unemployment stemmed primarily from inadequate education, lack of job training, and geographical remoteness (particularly among Alaska Native residents), not from an influx of out-of-state workers (\href{https://supreme.justia.com/cases/federal/us/437/518/}{437 U.S.\ at 526--27 \& n.10}). Moreover, even assuming nonresidents contributed to the problem by taking jobs before untrained residents completed training, Alaska Hire was grossly overbroad: it granted an identical preference to a highly skilled, fully employed Alaskan resident as to an unskilled, chronically unemployed resident enrolled in job training (\href{https://supreme.justia.com/cases/federal/us/437/518/}{437 U.S.\ at 527--28}).

Third, the Court rejected Alaska's defense under \textit{McCready v. Virginia}, 94 U.S.\ 391 (1877), that state ownership of oil and gas (acquired under the Alaska Statehood Law) exempted Alaska Hire from the Clause (\href{https://supreme.justia.com/cases/federal/us/437/518/}{437 U.S.\ at 528--31 \& n.11}). Citing \textit{Baldwin v. Montana Fish and Game Comm'n}, 436 U.S.\ 371, 385 (1978), Brennan emphasized that state resource ownership is not an absolute exception, but merely one factor in the analysis. Here, \S\ 38.40.050(a) swept in private employers with no contractual relationship with the State, no work on state land, and no state funding---reaching subcontractors, suppliers, refineries, and distribution networks caught in the ``economic ripple effect'' of oil and gas development (\href{https://supreme.justia.com/cases/federal/us/437/518/}{437 U.S.\ at 529--31}).

Fourth, Justice Brennan fortified this result by invoking the ``mutually reinforcing relationship'' between Article IV, \S\ 2, and the Commerce Clause, both descended from the Fourth Article of the Articles of Confederation (\href{https://supreme.justia.com/cases/federal/us/437/518/}{437 U.S.\ at 531--32 \& n.16}). Drawing on \textit{West v. Kansas Natural Gas Co.}, 221 U.S.\ 229 (1911), \textit{Pennsylvania v. West Virginia}, 262 U.S.\ 553 (1923), and \textit{Foster-Fountain Packing Co. v. Haydel}, 278 U.S.\ 1 (1928), Brennan stressed that Alaska's North Slope oil and gas were destined for interstate commerce and deemed by Congress to be of ``profound national importance'' under the Trans-Alaska Pipeline Authorization Act and the Alaska Natural Gas Transportation Act of 1976 (\href{https://supreme.justia.com/cases/federal/us/437/518/}{437 U.S.\ at 532--34 \& nn.17--18}). Permitting a State to hoard the employment benefits of nationally vital resources would invite interstate economic retaliation and undermine the constitutional premise that the States ``must sink or swim together'' (\href{https://supreme.justia.com/cases/federal/us/437/518/}{437 U.S.\ at 534}).

\subsection{Concurring Opinions}

None. All nine Justices joined Justice Brennan's opinion in full.

\subsection{Dissenting Opinions}

None.

\section{Examples: Future Applications}

\begin{enumerate}
    \item \textbf{Same-Side (Unconstitutional): State-Mandated Local Hiring Preference on Private Renewable-Energy Supply Chains.} Suppose Wyoming leases state trust land to private wind developers and enacts a statute requiring all wind-farm lessees, their contractors, turbine-component suppliers, and in-state transmission operators to hire qualified Wyoming residents before any nonresidents. Under \textit{Hicklin}, this law violates Article IV, \S\ 2: pursuing private employment is a fundamental common calling, Wyoming's land ownership cannot justify coercing downstream private employers and suppliers caught in the economic ripple effect, and a flat preference for all residents fails \textit{Toomer}'s tailoring requirement.

    \item \textbf{Same-Side (Unconstitutional): Blanket State Residency Requirement for Professional Licensure.} Suppose a State bars nonresidents from obtaining a license to practice law or work as an escrow officer within the State, asserting a generalized interest in local accountability and employment. Just as the Court held in \href{https://supreme.justia.com/cases/federal/us/470/274/}{\textit{Supreme Court of New Hampshire v. Piper}, 470 U.S.\ 274 (1985)}, and as the Fifth Circuit recently applied \textit{Hicklin} and \textit{Toomer} in \href{https://law.justia.com/cases/federal/appellate-courts/ca5/25-20524/25-20524-2026-08-04.html}{\textit{Polk v. Crawford}, No.\ 25-20524 (5th Cir.\ 2026)}, excluding nonresidents from a protected common calling without showing that nonresidents are a peculiar source of a distinct regulatory evil---and without using less discriminatory alternatives---violates the Privileges and Immunities Clause.

    \item \textbf{Opposite-Side (Constitutional): Differential Fees for Recreational Hunting or Non-Fundamental Benefits.} Suppose Montana charges nonresidents significantly higher license fees than residents for purely recreational big-game elk hunting, or a State charges out-of-state students higher tuition at public universities. Under the companion doctrine of \textit{Baldwin v. Montana Fish and Game Comm'n}, 436 U.S.\ 371 (1978), and related precedent (\href{https://www.law.cornell.edu/constitution-conan/article-4/section-2/clause-1/privileges-and-immunity-clause-doctrine-and-practice}{Cornell LII Constitution Annotated}), recreational sport and state-subsidized higher education do not involve a fundamental livelihood or ``common calling'' bearing on the vitality of the Nation as a single economic entity, so the Privileges and Immunities Clause does not bar reasonable resident preferences.

    \item \textbf{Opposite-Side (Constitutional): Direct State Employment in Core Political or Policymaking Offices.} Suppose a State requires that its own directly employed state troopers, cabinet officials, or civil-service policymakers be bona fide state residents. Unlike Alaska Hire's private-sector ripple-effect mandate, direct employment by the State itself in sovereign governance functions implicates the State's existence as a separate political community and its direct proprietary/sovereign capacity rather than private-market discrimination against nonresidents (\href{https://www.law.cornell.edu/constitution-conan/article-4/section-2/clause-1/privileges-and-immunity-clause-doctrine-and-practice}{Cornell LII Constitution Annotated}).

    \item \textbf{Fence-Sitter (Close / Fact-Dependent Case): Municipal Quota Requiring Resident Hiring on City-Funded Public Works Projects.} Suppose a financially distressed city enacts an ordinance requiring private contractors working directly on city-funded construction projects to ensure that 40\% of their workforce consists of city residents, citing severe urban unemployment and population flight. As the Supreme Court later confronted in \href{https://supreme.justia.com/cases/federal/us/465/208/}{\textit{United Building \& Construction Trades Council v. Mayor of Camden}, 465 U.S.\ 208 (1984)}, this sits squarely on the doctrinal fence: unlike \textit{Hicklin}, the city is directly expending its own public funds on its own public works rather than regulating attenuated private suppliers, so its proprietary interest is far stronger; yet because the ordinance still burdens private construction employment (a protected common calling), it is not exempt from Article IV, \S\ 2, and its constitutionality turns on a fact-intensive remand over whether nonresidents are a peculiar source of the city's economic blight and whether the 40\% quota is narrowly tailored.
\end{enumerate}

\section{Critique}

\textbf{Originalist and Textualist Perspectives:} Textualist and originalist scholars have debated both the breadth of Article IV, \S\ 2, and the Court's methodology in \textit{Hicklin}. On one hand, originalist defenders of federalism have criticized the Burger Court's line of Article IV decisions---including \textit{Hicklin} and \textit{Piper}---for deploying the Privileges and Immunities Clause to impose ``unfettered national markets for occupations'' at the expense of traditional state sovereignty over local economic regulation and state-owned resources (\href{https://www.heritage.org/the-constitution/report/the-originalism-revolution-turns-30-evaluating-its-impact-and-future}{Heritage Foundation, \textit{The Originalism Revolution Turns 30}}). On the other hand, strict textualists and structural originalists such as William Crosskey and John Hart Ely have criticized the Court's reliance on \textit{Corfield v. Coryell}'s ``fundamental rights'' filter (which \textit{Hicklin} inherited alongside \textit{Baldwin}) as an atextual gloss on a clause that simply commands equal treatment of out-of-state citizens across all state-conferred privileges, arguing that an anti-discrimination reading rooted in the Fourth Article of the Articles of Confederation supports \textit{Hicklin}'s outcome while questioning the Court's malleable balancing tests (\href{https://digitalcommons.law.uw.edu/cgi/viewcontent.cgi?article=2277&context=wlr}{Washington Law Review}; \href{https://constitution.congress.gov/browse/essay/artIV-S2-C1-1/ALDE_00013777/}{Constitution Annotated}).

\textbf{Progressive and Doctrinal Perspectives:} Doctrinal and progressive commentators have praised \textit{Hicklin} for reinforcing national economic unity and preventing resource-rich States from erecting beggar-thy-neighbor labor walls that disproportionately harm mobile workers (\href{https://scholarship.law.duke.edu/cgi/viewcontent.cgi?article=2677&context=dlj}{Duke Law Journal, \textit{Domicile Preferences in Employment}}). At the same time, doctrinal critics have pointed out two lingering tensions in Justice Brennan's opinion. First, by blurring Article IV Privileges and Immunities doctrine with dormant Commerce Clause resource cases---even though the dormant Commerce Clause contains a ``market participant'' exception that Article IV does not fully mirror, as later revealed in \href{https://supreme.justia.com/cases/federal/us/465/208/}{\textit{Camden}}---the Court left lower courts and legislatures uncertain about how much proprietary nexus is required to sustain a narrower local-hire law (\href{https://scholarship.law.duke.edu/cgi/viewcontent.cgi?article=1235&context=alr}{Alaska Law Review, \textit{Local Hire Laws}}). Second, from a social-welfare perspective, some commentators noted that \textit{Hicklin} left States like Alaska with few constitutionally viable tools to ensure that resource booms translate into jobs for historically marginalized, undertrained local populations (such as rural Alaska Natives) without triggering strict interstate-comity invalidation.

\textbf{Underlying Competing Values:} \textit{Hicklin v. Orbeck} exposes the enduring constitutional tension between \textbf{national economic union and interstate comity}---the principle that citizens of every State share a single national labor market free from protectionist retaliation---and \textbf{state sovereignty and local self-determination}---the desire of a State's political community to capture the economic benefits of its own publicly owned natural resources to relieve acute local unemployment and poverty.

\section{Key Quotations}

\begin{quote}
``[I]t does bar discrimination against citizens of other States where there is no substantial reason for the discrimination beyond the mere fact that they are citizens of other States.'' \ldots\ A ``substantial reason for the discrimination'' would not exist, the Court explained, ``unless there is something to indicate that noncitizens constitute a peculiar source of the evil at which the [discriminatory] statute is aimed.''
\end{quote}
\noindent---Justice Brennan, quoting \textit{Toomer v. Witsell}, 334 U.S.\ 385, 396, 398 (1948), in \href{https://supreme.justia.com/cases/federal/us/437/518/}{\textit{Hicklin v. Orbeck}, 437 U.S.\ 518, 525--26 (1978)} (\href{https://www.courtlistener.com/opinion/109911/hicklin-v-orbeck/}{CourtListener}).

\begin{quote}
``Alaska Hire simply grants all Alaskans, regardless of their employment status, education, or training, a flat employment preference for all jobs covered by the Act. A highly skilled and educated resident who has never been unemployed is entitled to precisely the same preferential treatment as the unskilled, habitually unemployed Arctic Eskimo enrolled in a job training program.''
\end{quote}
\noindent---Justice Brennan, \href{https://supreme.justia.com/cases/federal/us/437/518/}{\textit{Hicklin v. Orbeck}, 437 U.S.\ 518, 527--28 (1978)} (\href{https://www.courtlistener.com/opinion/109911/hicklin-v-orbeck/}{CourtListener}).

\begin{quote}
``Rather than placing a statute completely beyond the Clause, a State's ownership of the property with which the statute is concerned is a factor---although often the crucial factor---to be considered in evaluating whether the statute's discrimination against noncitizens violates the Clause.''
\end{quote}
\noindent---Justice Brennan, \href{https://supreme.justia.com/cases/federal/us/437/518/}{\textit{Hicklin v. Orbeck}, 437 U.S.\ 518, 529 (1978)} (\href{https://www.courtlistener.com/opinion/109911/hicklin-v-orbeck/}{CourtListener}).

\begin{quote}
``In sum, the Act is an attempt to force virtually all businesses that benefit in some way from the economic ripple effect of Alaska's decision to develop its oil and gas resources to bias their employment practices in favor of the State's residents. We believe that Alaska's ownership of the oil and gas that is the subject matter of Alaska Hire simply constitutes insufficient justification for the pervasive discrimination against nonresidents that the Act mandates.''
\end{quote}
\noindent---Justice Brennan, \href{https://supreme.justia.com/cases/federal/us/437/518/}{\textit{Hicklin v. Orbeck}, 437 U.S.\ 518, 531 (1978)} (\href{https://www.courtlistener.com/opinion/109911/hicklin-v-orbeck/}{CourtListener}).

\begin{quote}
``As Mr.\ Justice Cardozo observed in \textit{Baldwin v. G. A. F. Seelig, Inc.}, 294 U.\ S.\ 511, 523 (1935), the Constitution `was framed upon the theory that the peoples of the several states must sink or swim together, and that, in the long run, prosperity and salvation are in union, and not division.' ''
\end{quote}
\noindent---Justice Brennan, \href{https://supreme.justia.com/cases/federal/us/437/518/}{\textit{Hicklin v. Orbeck}, 437 U.S.\ 518, 534 (1978)} (\href{https://www.courtlistener.com/opinion/109911/hicklin-v-orbeck/}{CourtListener}).

\end{document}
